\section{A calibrated family of cheaper equality binders}
\label{sec:binders}

The machine's merged equality axiom uses a formula code as a fresh
variable label. Its particular tag is not logically essential. The following
result identifies exactly how much freedom the compiler has. This is a
new ordinary mathematical generalization of the inherited calibration; it
is not a newly compiled exact-table Lean theorem.

Write
\[
 P(a,b)=\frac{(a+b)(a+b+1)}2+b,\qquad
 I(a,b)=4P(a,b),\quad R(a,b)=4P(a,b)+1,
\]
\[
 Q_i(A)=4P(i,A)+2,\qquad E(a,b)=4P(a,b)+3,
 \qquad \bot=1,\quad \top=0.
\]
The exceptional identity $I(0,0)=0$ remains part of the numerical
calculus. All the arguments below are about literal natural-number codes.

\begin{theorem}[Equality binders fresh off the diagonal]
\label{thm:flexible-binder}
Let $f:\mathbb N^2\to\mathbb N$ be primitive recursive, and assume
\begin{equation}
 a\ne b\quad\Longrightarrow\quad
 f(a,b)\ne a\ \hbox{ and }\ f(a,b)\ne b.
 \label{eq:offdiag}
\end{equation}
Retain truth, numerical modus ponens, theorem-only generalization,
the six rows $\rho_1,\ldots,\rho_6$, the exact gated seed bundle,
and the dominated $\rho_8$ instances $U\le A$ of the accepted calculus.
Replace only its merged equality row by
\begin{equation}
 G_f(a,b,c)=
 E(c,a)\to\bigl(E(c,b)\to Q_{f(a,b)}E(a,b)\bigr).
 \label{eq:gf}
\end{equation}
Call the resulting numerical calculus $C_f$. Uniformly in an index for
$f$, ZFC proves that condition~\eqref{eq:offdiag} implies
\[
 \operatorname{Con}(C_f)\ \longleftrightarrow\
 \operatorname{Con}(\mathrm{ZFC}).
\]
In particular, whenever ZFC proves condition~\eqref{eq:offdiag} for a
chosen $f$, it proves the displayed equivalence without an additional
premise. This applies to every explicit polynomial binder listed below.
Here $\operatorname{Con}(C_f)$ means that native code $1$ is not derivable.
The assertion is a refutability calibration; it does not identify all
native theorem codes of $C_f$ with those of its parent.
\end{theorem}

\subsection{The equality derivations are independent of the binder}

The inherited finite derivation from dominated $\rho_8$ recovers the
unrestricted classical implicational axiom using only that row, truth, and
modus ponens. It therefore works without change in $C_f$. The retained $\rho_6$
then supplies explosion, as in \texttt{KCode222.c222\_efq}. In the following
proof, classical deduction, contraposition, contraction, and conjunction
manipulations abbreviate fixed Hilbert derivations from that basis.
Only closed theorems are generalized.

Put $\exists_rX=\neg Q_r\neg X$, with $\neg X=I(X,1)$. The retained
$\rho_2$ and closed generalization give universal monotonicity, and
contraposition gives existential monotonicity:
\[
 \vdash X\to Y\quad\Longrightarrow\quad
 \vdash Q_rX\to Q_rY,\qquad
 \vdash\exists_rX\to\exists_rY.
\]
They also give
\begin{equation}
 \vdash Q_r(X\to Y)\to(\exists_rX\to\exists_rY).
 \label{eq:all-ex-binder}
\end{equation}
These are the literal \texttt{all\_ex} and \texttt{ex\_mono} constructions
in \texttt{KEqCodes222.lean}.

\paragraph{Reflexivity.}
Fix $x$ and set $r=f(x,x)$. Contracting $G_f(x,x,r)$ gives
\[
 E(r,x)\to Q_rE(x,x).
\]
Existential monotonicity and $\rho_1(r,x)$ give
$\exists_rQ_rE(x,x)$. Applying $\rho_6(r,E(x,x))$ yields $E(x,x)$.
There is no freshness premise in this derivation: $r=x$ is permitted.

\paragraph{The helper.}
For $r=f(a,b)$, apply $G_f(a,b,a)$ to reflexivity $E(a,a)$. This gives
\begin{equation}
 \vdash E(a,b)\to Q_rE(a,b).
 \label{eq:f-helper}
\end{equation}

\paragraph{Specialization at the helper binder.}
The $K$ axiom and universal monotonicity give
\[
 Q_rE(a,b)\to Q_r\bigl(E(r,a)\to E(a,b)\bigr).
\]
Combine this with \eqref{eq:all-ex-binder} and $\rho_1(r,a)$ to derive
$Q_rE(a,b)\to\exists_rE(a,b)$. Existential monotonicity applied to
\eqref{eq:f-helper}, followed by $\rho_6$, now gives
\begin{equation}
 \vdash Q_rE(a,b)\to E(a,b).
 \label{eq:f-spec}
\end{equation}
Composing $G_f$ with \eqref{eq:f-spec} recovers the unrestricted equality
row
\[
 \rho_7(a,b,c)=E(c,a)\to\bigl(E(c,b)\to E(a,b)\bigr).
\]
Thus reflexivity, the helper, and $\rho_7$ are recovered for an
\emph{arbitrary} function $f$. The numerical freshness condition enters
only the next step and the semantic soundness argument.

\subsection{The diagonal exception in the completeness argument}

Let $\operatorname{Eind}(x,y)$ be the standard equality-free
indiscernibility formula
\[
 \bigl[\forall v\,(R(v,x)\leftrightarrow R(v,y))\bigr]
 \wedge
 \bigl[\forall v\,(R(x,v)\leftrightarrow R(y,v))\bigr],
\]
where the displayed $v$ is the particular fresh label $x+y+1$.
The connectives have the same implication/falsity expansion as
\texttt{KEind222.lean}.

The inherited equality-free vacuity, quantifier-free specialization,
and alpha-renaming derivations use $\rho_1$--$\rho_6$, recovered
$\rho_7$, equality reflexivity, membership vacuity, and the classical
basis. They do not require equality vacuity at every fresh label.
Membership vacuity remains a projection of the unchanged gated bundle.

We prove
\begin{equation}
 \vdash E(x,y)\to\operatorname{Eind}(x,y).
 \label{eq:eq-to-ind-f}
\end{equation}
For $x\ne y$, choose $w=f(x,y)$. Condition \eqref{eq:offdiag} makes
$w$ different from both $x$ and $y$. The retained membership-substitution
rows and equality symmetry give implications from $E(x,y)$ to both
membership biconditionals with auxiliary label $w$. The helper
\eqref{eq:f-helper} gives $E(x,y)\to Q_wE(x,y)$; universal monotonicity
therefore puts both biconditionals under $Q_w$. Conjunction introduction
and the existing equality-free alpha-renaming proof change $w$ to
$x+y+1$, proving \eqref{eq:eq-to-ind-f}. This is the literal
\texttt{KEind222.Eq\_Eind} proof with $f$ in place of $R$.

For $x=y$, each biconditional $R(v,x)\leftrightarrow R(v,x)$ and
$R(x,v)\leftrightarrow R(x,v)$ is a propositional theorem. Closed
generalization and conjunction introduction give
$\operatorname{Eind}(x,x)$ directly, and $K$ gives
$E(x,x)\to\operatorname{Eind}(x,x)$. No freshness of $f(x,x)$ is
needed. This explicit diagonal branch is what permits the compact
choice $f(0,0)=0$.

Existential monotonicity and $\rho_1$ applied to
\eqref{eq:eq-to-ind-f} give $\exists_i\operatorname{Eind}(i,j)$.
The other translated equality rules---Euclid, membership substitution,
and arbitrary fresh vacuity of $\operatorname{Eind}$---use only the
unchanged equality-free quantifier arguments. Consequently every finite
ordinary Tarski derivation translates after replacing equality by
$\operatorname{Eind}$, including a refutation translated to native code
$1$.

\paragraph{Dependency audit.}
In the inspected \texttt{KThmK.lean}, the helper appears only in its rule
definition and the proof of \texttt{eq\_refl}. In
\texttt{KEind222.lean}, its only proof use is
\texttt{Eq\_Eind}; the other occurrence is the compiler's direct helper
case. The \texttt{eq\_refl} use is replaced by the arbitrary-$f$ reflexivity
derivation above. The \texttt{Eq\_Eind} use is replaced by the explicit
off-diagonal and diagonal arguments above. The compiler's helper case
uses the directly derived formula~\eqref{eq:f-helper}. In particular,
the alpha-renaming argument does not hide another fresh-helper
assumption.

\subsection{Transfer of the two consistency implications}

For the lower implication, repeat the exact interface of
\texttt{KEind222.c222\_complete} with $C_f$ in place of $C222$.
Let $T$ consist of the named formulas whose equality interpretations have
derivable native codes. If $C_f$ has no refutation, first-order
completeness and the interpreted Tarski proof yield a nonempty
irreflexive model of $T$.

For every equality-free named matrix $\psi$, its equality interpretation
is unchanged. Double-negation padding gives an equivalent matrix with
code greater than the role labels $0,1$. The exact gated seed bundle at
that code remains an axiom of $C_f$. Hence the entire argument of
\texttt{KSeedModel222.seed\_eqFree}, including its native-code/private-label
coupling and low-code treatment, remains valid. The model has the seed
property for every equality-free relation.
\texttt{KIndisc222.fullSeed\_quotient} takes the indiscernibility quotient
and obtains the seed property for every named relation, equality included.
The accepted lower seed-theory calibration then supplies
$\operatorname{Con}(C_f)\to\operatorname{Con}(\mathrm{ZFC})$.

This route deliberately does not substitute inside an arbitrary old
numerical bundle: that would change both its matrix code and its private
label. It uses the same exact equality-free bundles and semantic quotient
as the accepted proof.

For the upper implication, keep the original constructible-model
interpretation of the exact seed bundles in
\texttt{KUpper222.bundles\_in\_L}. All retained logical rows are unchanged.
If $a\ne b$, condition \eqref{eq:offdiag} means that changing the binder
of $Q_{f(a,b)}E(a,b)$ changes neither argument of the equality, so the two
antecedent equalities imply its consequent. If $a=b$, the quantified
consequent is reflexivity under every assignment, even when $f(a,a)=a$.
Thus $G_f$ is valid in every ordinary structure. Numerical decoding,
including the zero alias, and induction on finite derivations give
$\operatorname{Con}(\mathrm{ZFC})\to\operatorname{Con}(C_f)$.
This completes the proof of Theorem~\ref{thm:flexible-binder}.

\subsection{Sharpness and cheap arithmetic choices}

Condition \eqref{eq:offdiag} exactly characterizes validity of all $G_f$
rows in ordinary structures. Suppose $a\ne b$ but $f(a,b)$ equals one of
$a,b$. In a two-element structure, assign the initial value $0$ to every
relevant label, including $c$. Both antecedent equalities are true. In the
quantified consequent, assign value $1$ to the binder. Exactly one of the
distinct labels $a,b$ changes, so the equality is false. Thus the condition
is necessary as well as sufficient.

Let $s=a+b$. Then $P(a,b)\ge s$. For every fixed integer $k\ge2$,
\[
 s>0\quad\Longrightarrow\quad kP(a,b)\ge ks>s\ge\max(a,b).
\]
The only exceptional pair is $(0,0)$, which lies on the permitted diagonal.
Also $P(a,b)+1>\max(a,b)$ for all $a,b$. Therefore the following binders
are calibrated:
\[
 2P(a,b),\qquad 4P(a,b),\qquad P(a,b)+1,\qquad
 4P(a,b)+j\quad(j\ge1),\qquad a+b+1.
\]
The unshifted pairing $P$ is excluded: $P(1,0)=1$ collides with its first
argument off the diagonal.

The accompanying arithmetic verifier checks exact rational-polynomial
identities and nonnegative coefficients after the exhaustive cases
$s=0$ and $s=t+1$. It does not infer the all-natural statement from
finite numerical samples. A separate finite equality-pattern audit has
a finite-support proof: the validity of $G_f$ depends only on the equality
pattern of the four labels $a,b,c,f(a,b)$ and at most four domain values,
including the quantified value. Its exhaustive cases independently check
the sharp structural condition above.

\paragraph{Additional harmless orientation changes.}
The two equality antecedents may be exchanged by classical implication
exchange. Simultaneously exchanging the two arguments in the binder
and the final equality is a global $a,b$ permutation of the schema,
followed if necessary by that same antecedent exchange. These alternatives
have the same calibrated refutations and may expose different repeated
source subtrees. Exchanging only the equality body while leaving an
asymmetric binder unchanged is not justified by this observation.

\subsection{Source opportunity and remaining boundary}

In the inspected K215 source, the complete merged-equality emitter occupies
primitive ordinals $[479,888)$; its $R$-binder constructor occupies
$[695,723)$. Its macro word is
\begin{quote}\small\ttfamily
168,166,182,168,121,182,166,121,181,166,121,182,180,119,119.
\end{quote}
The constructors $119,181,180,182$ mean $I,R,Q,E$, respectively, while
$166,121,168$ load $a,b,c$. Replacing the sole $181$ by $119$ produces
$G_{4P}$ and removes the relation-tag increment. Other binders need their
own literal source construction. Shorter arithmetic does not by itself
reduce the shared source DAG or physical state count; source association,
allocation, complete runtime, and blank initialization remain separate
obligations.

Removing the seed-bundle freshness comparisons is a different unresolved
problem. Restricting bundle emission to $a,b<c$ retains enough padded
matrices for the seed extraction, but restricts the simultaneously emitted
membership-vacuity schema. The existing completeness proof needs that
schema at arbitrary fresh labels. Its alpha-renaming theorem already uses
membership vacuity, so it cannot silently repair this restriction.

\paragraph{Pinned inherited sources.}
All cited modules and the literal source measurement refer to ProveIt
revision \texttt{2762b8a183830b1f5a30cf059ffb4d92e77cddde}.
The modules are under
\texttt{Computability/BusyBeaver/Lean/BusyBeaver/ZFC/}; the machine source
is \texttt{Research/ZFC/candidate-k215/sites215.json} relative to
\texttt{Computability/BusyBeaver/}. The arbitrary-function binder theorem,
the off-diagonal refinement, its sharpness argument, and the arithmetic
certificate are additions of the present investigation. The seed
calibration and the existing named/numerical proof infrastructure retain
their original attribution.
